\begin{figure}[H]
    \centering
\begin{tikzpicture}[
    scale=1.5,
    cota/.style={<->, >=stealth, thin, gray!80!black},  % Estilo das linhas de cota
    retangulo/.style={draw=gray!80!black, thick},        % Estilo do retângulo
    chamada/.style={gray!60, thin, dashed},              % Linhas de chamada/extensão
    eixo/.style={->, >=stealth, thick, gray!80!black}   % Estilo dos eixos x e y
]

    % --- Parâmetros de Dimensão ---
    \def\wE{2} % Largura do retângulo esquerdo
    \def\hE{1} % Altura do retângulo esquerdo
    
    \def\wD{1} % Largura do retângulo direito (menor)
    \def\hD{2} % Altura do retângulo direito (menor)
    
    % --- Coordenadas ---
    % Topo comum em y = \hE
    \coordinate (A_E) at (0, 0);                 % Canto inf. esq. do esquerdo
    \coordinate (B_E) at (\wE, 0);               % Canto inf. dir. do esquerdo
    \coordinate (C_E) at (\wE, \hE);             % Canto sup. dir. do esquerdo
    \coordinate (D_E) at (0, \hE);               % Canto sup. esq. do esquerdo

    \coordinate (A_D) at (\wE, \hE - \hD);       % Canto inf. esq. do direito
    \coordinate (B_D) at (\wE + \wD, \hE - \hD); % Canto inf. dir. do direito
    \coordinate (C_D) at (\wE + \wD, \hE);        % Canto sup. dir. do direito

    % --- 1. Desenho dos Retângulos (Alinhados no topo em y = \hE) ---
    \draw[retangulo] (A_E) rectangle (C_E); % Retângulo da esquerda
    \draw[retangulo] (A_D) rectangle (C_D); % Retângulo da direita
    \draw[retangulo, black, line width=3pt] (C_E) rectangle (B_E); 
    %\draw[retangulo, black, line width=3pt] (A_E) rectangle (B_E); 

    % --- 2. Números nos centros dos retângulos ---
    \node[font=\large\bfseries, gray!80!black] at ($(A_E)!0.5!(C_E)$) {L1};
    \node[font=\large\bfseries, gray!80!black] at ($(A_D)!0.5!(C_D)$) {L2};

    % --- 2. Eixos de Coordenadas x e y (Canto inf. esq. do retângulo da esquerda) ---
    %\draw[eixo] (0.2, 0.2) -- (0.8, 0.2) node[above, font=\small\bfseries] {$x$};
    %\draw[eixo] (0.2, 0.2) -- (0.2, 0.8) node[right, font=\small\bfseries] {$y$};

    % --- 3. Cotas Horizontais Inferiores (Eixo único com comprimentos separados) ---
    \def\offsetH{1.5} % Distância do nível inferior do retângulo da esquerda
    \def\yCota{-\offsetH}

    % Linhas de chamada verticais estendidas até a cota
    %\draw[chamada] (A_E) -- (0, \yCota - 0.15);
    %\draw[chamada] (B_E) -- (\wE, \yCota - 0.15);
    %\draw[chamada] (B_D) -- (\wE + \wD, \yCota - 0.15);

    % Cota do Retângulo da Esquerda (b1)
    \draw[cota] (0, \yCota) -- (\wE, \yCota) 
        node[midway, below=3pt, font=\small] {$476 \ \text{ cm}$};

    % Cota do Retângulo da Direita (b2)
    \draw[cota] (\wE, \yCota) -- (\wE + \wD, \yCota) 
        node[midway, below=3pt, font=\small] {$446 \ \text{ cm}$};

\end{tikzpicture}
    %\caption{Dois retângulos alinhados no topo com cotas separadas na base}
    %\label{fig:retangulos_alinhados}
\end{figure}